% generated by src/feasibility.py -- do not edit
\begin{figure}
  \centering
  \includegraphics[width=\linewidth]{demand_feasibility.pdf}
  \caption{Why the model has no schedule for a summer day. Each panel is one day's hourly demand (100~MW of new solar, baseline data). Lignite output is held constant over the day, so the remaining plants -- hydro, gas and solar -- must follow the whole of the daily range of demand, and between their minimum and maximum outputs they span at most 1,802~MW after losses. The grey band is what is reachable: what they can cover on top of the best flat lignite level. The winter day fits inside it. The summer day has a deeper night-time trough relative to its evening peak, and no lignite level fits: set low enough for the trough, the band cannot reach the peak, and vice versa. Marked hours fall outside the band at that level.}
  \label{fig:demand_feasibility}
\end{figure}

\begin{table}
  \centering
  \caption{Feasibility of a winter and a summer demand day. Range is peak minus minimum hourly demand. Shortfall is how far the daily range exceeds the flexible range -- what hydro, gas and solar can follow with lignite held flat -- and ``--'' means a feasible schedule exists. The flexible range grows with new solar because run.py does not cap solar by an hourly production profile; even so, the summer day does not become feasible. The verdicts depend only on plant output bounds, losses and demand, not on prices or emission factors, and are identical under the 2026 data. Each was confirmed by solving the LP.}
  \label{tab:demand_feasibility}
  % generated by src/feasibility.py -- do not edit
% shortfall: MW by which the daily range exceeds what
% hydro, gas and solar can follow with lignite flat; -- = feasible
\begin{tabular}{llrrrrrrr}
\toprule
 & & \multicolumn{3}{c}{Demand (MW)} & \multicolumn{4}{c}{Shortfall (MW), by new solar (MW)} \\
\cmidrule(lr){3-5}\cmidrule(lr){6-9}
Day & Season & Min & Peak & Range & 0 & 10 & 50 & 100 \\
\midrule
19 Jan 2021 & winter & 4{,}205 & 5{,}188 & 983 & -- & -- & -- & -- \\
3 Aug 2026 & summer & 2{,}796 & 4{,}645 & 1{,}849 & 147 & 137 & 97 & 47 \\
\midrule
\multicolumn{5}{l}{Flexible range (MW)} & 1{,}702 & 1{,}712 & 1{,}752 & 1{,}802 \\
\bottomrule
\end{tabular}

\end{table}
